% ==============================================================================
% CAPÍTULO 12: SOFTWARE COMO PRODUCTO (Orquestador del Capítulo)
% ==============================================================================
\preportadacapitulo{12}{Software como Producto}

\newpage
\section{Capítulo 12: Software como Producto}\label{cap:producto}

Este capítulo reúne los materiales de ayuda al usuario del sistema de la Beca Comedor: el manual de usuario y, como material previsto, un asistente conversacional y tutoriales guiados.

\subsection{12.1. Para el Usuario}

\subsubsection{12.1.1. Manual de usuario}

\paragraph{12.1.1.1. Alcance del manual.}

El manual describe las funciones del sistema que están implementadas a la fecha: 13 de los 26 casos de uso (CU-01, CU-04, CU-05, CU-06, CU-07, CU-08, CU-12, CU-14, CU-15, CU-19, CU-20, CU-21 y CU-22), junto con el servicio de avisos que los acompaña. Las funciones de los casos de uso restantes (entre ellas el envío de la postulación, la precalificación con inteligencia artificial y el control de asistencia) se agregarán al manual cuando se implementen. El sistema se usa desde un navegador web (panel web) y desde una aplicación móvil Android; la aplicación móvil cubre hoy el acceso, el perfil, los avisos y la consulta de las postulaciones del postulante. Las capturas de pantalla de cada pantalla descrita se incorporarán cuando el equipo las entregue.
% TODO: insertar capturas de pantalla de cada pantalla descrita (a cargo de Mau).

\paragraph{12.1.1.2. Perfiles de usuario.}

Los usuarios del sistema pertenecen a cuatro perfiles, y cada uno ve solo las opciones que le corresponden: el Postulante (estudiante de la UAGRM), el Administrador de la DUBSS, el Directivo y el Personal del Comedor. Al iniciar sesión el sistema lleva a cada persona al módulo de su perfil. La Tabla~\ref{tab:menu_por_rol} resume las opciones del panel web por perfil.

\begin{longtable}{>{\raggedright\arraybackslash}p{3.4cm} >{\raggedright\arraybackslash}p{9.4cm}}
\caption{Opciones del panel web según el perfil del usuario}\label{tab:menu_por_rol}\\
\toprule
Perfil & Opciones del menú \\
\midrule
\endfirsthead
\toprule
Perfil & Opciones del menú \\
\midrule
\endhead
\bottomrule
\endfoot
Sin sesión & Inicio, Iniciar sesión, Crear cuenta y Recuperar contraseña. \\
Postulante & Mis postulaciones, Resultados, Perfil y bandeja de avisos. \\
Administrador de la DUBSS & Dashboard, Lista priorizada, Finalistas, Auditoría y adjudicación, Publicar resultados, Actas, Administradores DUBSS, Directivos, Personal del Comedor, Perfil y bandeja de avisos. \\
Directivo & Dashboard de indicadores. \\
\end{longtable}
\noindent\textit{Nota.} Elaboración propia a partir de las opciones de menú del panel web. Las opciones heredadas del prototipo inicial (Menús, Pedidos, Mi Menú y Mis Pedidos) no forman parte de este producto.

\paragraph{12.1.1.3. Guía para el Postulante.}

\textbf{Crear una cuenta y entrar (CU-01).} El Postulante abre la pantalla \textit{Crear cuenta} (en el panel web o en la aplicación móvil), completa sus nombres, apellidos, carné de identidad, registro universitario (9 dígitos), carrera, correo electrónico y contraseña, y confirma. El sistema crea la cuenta con el perfil de Postulante y abre la sesión. El carné, el registro universitario y el correo no pueden estar repetidos. Para entrar en otras ocasiones, usa \textit{Iniciar sesión} con su correo y contraseña. Tras cinco intentos fallidos la cuenta se bloquea durante un tiempo, aunque luego se escriba la contraseña correcta. Para salir usa la opción \textit{Cerrar sesión} del menú de usuario.

\textbf{Recuperar la contraseña (CU-15).} En la pantalla de acceso, el Postulante pulsa \textit{¿Olvidó su contraseña?}, escribe su correo y recibe siempre el mismo mensaje, exista o no la cuenta. Si existe, le llega un enlace de un solo uso que vence en 60 minutos; al abrirlo define una contraseña nueva y se cierran sus otras sesiones.

\textbf{Gestionar el perfil (CU-15).} En \textit{Perfil} el Postulante ve sus datos, edita los de contacto (el carné de identidad no se puede cambiar) y puede cambiar su contraseña escribiendo primero la actual.

\textbf{Consultar el estado de una postulación (CU-04).} En \textit{Mis postulaciones} el Postulante ve la lista de sus postulaciones, de la más reciente a la más antigua, con la gestión académica, la fecha de envío y el estado. Puede filtrar por gestión. Al abrir una postulación ve el detalle con una línea de tiempo con la fecha y la hora de cada cambio y los pasos que aún faltan. Según el estado, se agrega información: en \textit{devuelta}, las observaciones del personal y la opción de corregir y reenviar; en \textit{archivada} o \textit{no adjudicada} por rechazo, el motivo registrado; en el caso de un finalista, el plazo y los documentos que debe entregar en la DUBSS. Hasta que la DUBSS publique los resultados, una postulación ya decidida se sigue mostrando como \textit{calificada}. La misma consulta está disponible en la aplicación móvil.

\textbf{Ver los resultados (CU-22).} Una vez publicados los resultados de la convocatoria, el Postulante los consulta en \textit{Resultados}. Cada postulante ve su propio resultado.

\textbf{Avisos.} El sistema avisa al Postulante, en la bandeja de avisos, con una notificación en el celular y por correo, cuando es notificado como finalista, cuando se registra el resultado de la auditoría de su expediente, cuando su postulación es devuelta, archivada o rechazada, cuando se publican los resultados y cuando la DUBSS emite un comunicado. Al tocar un aviso se abre la postulación correspondiente.

\paragraph{12.1.1.4. Guía para el Administrador de la DUBSS.}

\textbf{Gestionar usuarios y roles (CU-14).} En las opciones \textit{Administradores DUBSS}, \textit{Directivos} y \textit{Personal del Comedor} el Administrador ve las cuentas del personal por perfil y puede crear una cuenta, editarla, cambiar su rol, desactivarla y reactivarla. Al crear una cuenta el sistema genera una contraseña temporal que se muestra una sola vez. Las cuentas nunca se borran; no es posible desactivar la propia cuenta ni la del último administrador activo.

\textbf{Visualizar la lista priorizada (CU-06).} En \textit{Lista priorizada} el Administrador elige la convocatoria (por defecto la más reciente) y ve las postulaciones calificadas ordenadas de mayor a menor puntaje de vulnerabilidad, con su posición, nombre, carné, carrera y puntaje. Puede buscar por nombre o carné, filtrar por carrera y exportar la lista filtrada a PDF; la exportación queda anotada en la bitácora. La lista siempre muestra el puntaje vigente.

\textbf{Notificar finalistas, registrar la entrega y publicar comunicados (CU-21).} En \textit{Finalistas} el Administrador realiza tres tareas: (a) notificar a los finalistas elegidos indicando el plazo, el lugar, el horario y los documentos a presentar; (b) marcar la documentación recibida, con la fecha y las observaciones, cuando el finalista entrega sus papeles; y (c) redactar un comunicado con título y contenido, elegir los destinatarios (todos o solo los finalistas), revisar la vista previa y publicarlo. Sin la documentación recibida no es posible auditar el expediente.

\textbf{Auditar el expediente de un finalista (CU-07).} Desde la fila del finalista, con la opción \textit{Auditar}, el Administrador abre el expediente: datos del postulante, respuestas socioeconómicas, documentos (recibidos o faltantes, que se abren con un acceso temporal) y puntaje de la inteligencia artificial. Registra el resultado de la auditoría (\textit{Validado}, \textit{Observado} o \textit{No válido}), las observaciones (hasta 500 caracteres) y el puntaje académico de 0 a 30, obligatorio cuando el resultado es \textit{Validado}. El sistema muestra el puntaje total (máximo 100), suma del puntaje de la inteligencia artificial (hasta 70) y del puntaje académico (hasta 30), y avisa al postulante del resultado. Un finalista puede tener varias auditorías y vale la última.

\textbf{Devolver, archivar o rechazar desde el expediente (CU-19 y CU-20).} Desde el expediente, el Administrador puede devolver la postulación para corrección (CU-19), con un motivo obligatorio de hasta 500 caracteres; la opción solo se habilita si la postulación está calificada y fue auditada después de la última devolución. También puede archivarla o rechazarla (CU-20), con motivo, cuando no procede y no tiene arreglo. Estas dos acciones son definitivas. En los tres casos se avisa al postulante con el motivo.

\textbf{Auditar y adjudicar becas (CU-05).} En \textit{Auditoría y adjudicación} el Administrador elige la convocatoria y ve el panel con los cupos totales, los adjudicados, los no adjudicados, los cupos disponibles y el puntaje mínimo. Selecciona como finalistas a los postulantes mejor puntuados (el sistema propone la cantidad configurada en la convocatoria, que puede ajustar), y decide desde el expediente de cada finalista auditado con \textit{Adjudicar} o \textit{No adjudicar}. Para adjudicar, la última auditoría debe ser \textit{Validado}, el puntaje total debe alcanzar el mínimo y deben quedar cupos; para no adjudicar se exige un motivo. La decisión es final y el postulante no recibe aviso hasta la publicación de los resultados.

\textbf{Publicar resultados (CU-22).} En \textit{Publicar resultados} el Administrador elige la convocatoria, revisa el resumen de adjudicados y no adjudicados y pulsa \textit{Publicar}. El sistema advierte que la publicación no se puede deshacer, y solo la permite cuando no quedan finalistas pendientes de decisión; al publicar avisa a cada finalista de su resultado.

\textbf{Generar el acta de adjudicación (CU-08).} En \textit{Actas} el Administrador ve las actas generadas y las convocatorias con resultados publicados que aún no tienen acta. Elige una convocatoria, puede ver una vista previa marcada como borrador y pulsa \textit{Generar PDF}; el sistema advierte que el acta es definitiva, la genera con los adjudicados (nombre, carné, carrera y puntaje final) y la guarda. El acta se descarga o imprime desde la misma pantalla con \textit{Descargar}; la descarga queda anotada en la bitácora.

\paragraph{12.1.1.5. Guía para el Directivo.}

\textbf{Visualizar el dashboard de indicadores (CU-12).} Al iniciar sesión el Directivo abre \textit{Dashboard}, que muestra cuatro indicadores y los gráficos de postulaciones y de asistencia por mes de la gestión académica vigente, junto con la fecha de la última actualización. Puede elegir otra gestión o una convocatoria, pulsar \textit{Aplicar} para actualizar la vista y, si lo desea, pulsar \textit{Descargar PDF} para obtener el informe de lo que está viendo. Si no hay datos, el sistema lo informa y muestra los indicadores en cero.

\paragraph{12.1.1.6. Mensajes y problemas frecuentes.}

\begin{longtable}{>{\raggedright\arraybackslash}p{5.2cm} >{\raggedright\arraybackslash}p{7.6cm}}
\caption{Situaciones frecuentes y qué hacer}\label{tab:problemas_frecuentes}\\
\toprule
Situación & Qué hacer \\
\midrule
\endfirsthead
\toprule
Situación & Qué hacer \\
\midrule
\endhead
\bottomrule
\endfoot
La cuenta está bloqueada por intentos fallidos & Esperar a que termine el bloqueo, que aumenta si se sigue fallando, o recuperar la contraseña. \\
El enlace de recuperación no funciona & Es de un solo uso y vence a los 60 minutos; solicitar uno nuevo. \\
No se habilita \textit{Auditar} & Registrar primero la documentación recibida del finalista. \\
No se puede devolver una postulación & Debe estar calificada y haber sido auditada después de la última devolución. \\
No se puede publicar los resultados & Quedan finalistas sin decisión; el mensaje indica cuántos faltan. \\
Dos personas ejecutan la misma acción a la vez & La segunda recibe un aviso de conflicto y no se repite la operación. \\
\end{longtable}
\noindent\textit{Nota.} Elaboración propia a partir de las reglas de los casos de uso.

\subsubsection{12.1.2. Chatbots}

Un asistente conversacional (\textit{chatbot}) es un programa que responde preguntas en lenguaje natural. El sistema todavía no incluye uno; esta sección describe lo previsto. El asistente resolvería las dudas frecuentes de los postulantes y del personal (requisitos de la beca, documentos a presentar, significado de cada estado de la postulación, plazos y cómo corregir una postulación devuelta) con respuestas basadas en el manual de usuario y en la información oficial de la convocatoria. Estaría disponible dentro del panel web y de la aplicación móvil, y derivaría a la DUBSS las consultas que no pueda resolver. Las consultas sobre la situación personal de un postulante se responderían solo con el estado que el propio usuario autenticado puede ver. La tecnología, las fuentes de conocimiento y el alcance del asistente aún no están definidos; se acordarán con la DUBSS.
% TODO: definición del chatbot a cargo de Mau.

\subsubsection{12.1.3. Tutoriales}

Todavía no existen tutoriales. Está previsto incorporar un recorrido guiado (\textit{tour}) dentro del propio software, que oriente al usuario la primera vez que ingresa y pueda repetirse desde el menú de ayuda. El recorrido resaltaría paso a paso los elementos de la pantalla según el perfil: para el Postulante, cómo completar y enviar la postulación, cargar los documentos y consultar el estado; para el personal de la DUBSS, cómo revisar la lista priorizada, auditar un expediente y publicar resultados. Se complementaría con videos cortos y fichas de ayuda por caso de uso, que se elaborarían cuando las funciones respectivas estén implementadas. La herramienta del recorrido guiado y el calendario de elaboración se definirán más adelante.
